% ============================================================================
% Appendix G — Architecture positioning: frequently asked questions
% Status: COMPLETE.
%
% Scope: this appendix answers the comparative-positioning question that
% reduces the entire architecture to its simplest skeptical form: "if we
% are willing to spend an exact-match-labelled batch every cycle anyway,
% what does CAEM contribute that plain fine-tuning on the same batch
% would not?". The body of the thesis carries the full architectural
% specification in formal academic register; this appendix lays out the
% six load-bearing differentiators that justify the architecture even
% under that simplest skeptical framing.
% ============================================================================

\chapter[Architecture positioning]{Architecture positioning, frequently asked questions}
\label{app:positioning-qa}

The body chapters specify what CAEM does and how. This appendix collects the comparative-positioning and design-justification questions a reviewer in an adjacent area would ask, and answers each by naming the specific mechanism CAEM contributes that the alternative would lack. The first part answers \emph{why CAEM rather than the simpler alternative of plain exact-match-labelled fine-tuning at the same labelling budget} along six differentiators. The second part answers \emph{why the calibration discipline and cycle-boundary ordering chosen, rather than the straightforward alternatives}.

\part*{Part I: Positioning against plain fine-tuning}

% ---------------------------------------------------------------------------- %

\needspace{4\baselineskip}
\section[Per-query precision contract]{What does the per-query precision contract give a user that plain fine-tuning cannot?}
\label{app:qa-precision-contract}

Plain fine-tuning at the same labelling budget produces a model with an average exact-match accuracy on the held-out evaluation set; the deployed system serves whatever the model generates, with no per-query confidence axis. CAEM's calibrated-probability composite paired with the fixed-threshold storage gate produces a different deliverable: at every query the system outputs a calibrated probability $\ustored$ that the served answer is correct, and at the storage threshold $\tau_{\text{store}} = 0.60$ the calibration fold's empirical precision carries over to the stored answers when new questions resemble the calibration fold (\Cref{thm:purity}; the assumption held on the training panel and failed on TruthfulQA); the per-cycle composite refit absorbs distribution drift into the signal-to-probability mapping so that the same threshold continues to address the same precision target as the generator improves.

The user-facing consequence is the status tag and the abstain class. Every answer is shown with its probability and a tag, and a query whose composite is below the deferral threshold and whose grounding is below the abstention floor returns an insufficient-evidence message instead of a guessed answer. On the evaluation fold the tag predicts correctness: answers tagged verified were right $0.65$ to $0.75$ of the time at every cycle, against $0.51$ for B1 over all questions (\Cref{tab:status-tags}). A plain fine-tune serves every answer with no such signal; abstention needs a separate uncertainty estimate, such as the semantic-entropy baseline B8, which withholds $10.5$ percent of answers. For deployments where confident but wrong answers cost more than refusals (clinical question answering, legal document review, customer-facing knowledge serving), this per-answer signal is the deliverable that matters. Chatbots have been documented fabricating every reference offered for a clinical explanation \cite{alkaissi2023artificial}, up to 91.4\% of chatbot-generated references in systematic-review searches are hallucinated \cite{info:doi/10.2196/53164}, and TruthfulQA's inverse-scaling result \cite{lin-etal-2022-truthfulqa} shows that scale alone does not close this regime.

% ---------------------------------------------------------------------------- %

\needspace{4\baselineskip}
\section[Memory-direct tier at scale]{The memory-direct tier saves inference cost; does that saving survive at small deployment scale?}
\label{app:qa-memory-direct}

A memory-direct query serves an answer through a sentence-embedding cosine search over the FAISS-backed episodic memory and returns the stored answer at retrieval-only cost. A zero-shot generation through the same generator pays the full forward-pass cost. A retrieval-augmented generation pays the FAISS lookup cost plus the cross-encoder reranker cost plus the forward-pass cost over a longer context. The time per tier could not be separated in this run (\Cref{sec:tier-cost}), but a memory-direct answer skips generation and retrieval, the two most expensive steps of the other tiers.

The saving depends on how often questions repeat. In this study every evaluation question was new, so the memory-direct tier answered $7$ of $9{,}000$ evaluation queries over six cycles; on a probe that re-asked stored questions it answered $42$ of $50$ exact repeats and $19$ of $50$ paraphrases, $41$ and $19$ of them correctly (\Cref{tab:tier1-probe}). At the deployment scale projected in \Cref{sec:econ-deployment} the saving grows with traffic, but only in proportion to the share of repeated questions, which this study does not measure. Plain fine-tuning has no equivalent mechanism: every query pays the full generation cost however often the question has been seen, because the knowledge sits in the weights rather than in a separable episode.

% ---------------------------------------------------------------------------- %

\needspace{4\baselineskip}
\section[Asymmetric rollback vs manual backup]{What does asymmetric rollback give the operator that a manual checkpoint backup does not?}
\label{app:qa-asymmetric-rollback}

The retention guard at every cycle boundary measures the post-fine-tune model on a three-probe multi-modal panel (MMLU, TriviaQA-test, CommonsenseQA-test) and computes the per-probe retention ratio $\rho_{k} = \mathrm{score}_{k,\text{post}} / \mathrm{score}_{k,\text{pre}}$ for each probe $k$. When any worst-probe ratio falls below the registered floor of $0.93$ the cycle is declared aborted: the weights revert to the previous cycle's parameters, and the new memory written during the cycle is preserved unchanged. The architectural property is that the model is replaceable but the memory is not.

Plain fine-tuning has no equivalent guarantee even when the operator manually keeps the previous checkpoint. The reason is that plain fine-tuning at any cycle commits both the weights and the gradient signal that produced them; rolling the weights back rolls back the cycle's accumulated training experience entirely. CAEM's separation of weights from memory means a failed cycle costs only the cycle's fine-tune compute but does not lose the verified episodes that the cycle accumulated. The next cycle inherits the previous cycle's weights together with the larger memory, and the retroactive re-verification pass at the next cycle boundary refreshes the stored composites under the rolled-back generator. This asymmetric-rollback property is the architectural anchor for the memory-growth invariant established in \Cref{cor:stochastic-equilibrium}: memory accumulates strictly across cycles even when weights do not. A weights-only architecture has no monotonicity property to anchor.

% ---------------------------------------------------------------------------- %

\needspace{4\baselineskip}
\section[Label efficiency]{If both architectures need labelled data, how many fewer labels does CAEM use?}
\label{app:qa-sample-efficiency}

Plain fine-tuning uses a gold label for every training example, because the label is the training target. CAEM uses labels for a different purpose: the labelled calibration fold fits the verifier's composite and checks the precision at the storage cut-point, and the answers the generator trains on are labelled by the verifier. In the research run CAEM used about $2{,}100$ gold labels: the $1{,}500$ calibration questions, re-scored at every cycle, and the $603$ cold-start answers checked once (\Cref{sec:memory-store}). The vanilla fine-tuning baseline (B6) used a gold answer for each of its $23{,}500$ training questions, about eleven times as many.

With that difference, CAEM matched B6's pooled exact match ($0.535$ at C2 and $0.508$ at C5, against $0.512$), while B6 reached a lower composite hallucination metric ($0.103$ against $0.134$ at C5) and was the most accurate system on FEVER and CommonsenseQA (\Cref{tab:baseline-pooled}). The label saving is therefore real but not free: the verifier's labels are noisier than gold labels, about one stored answer in four is wrong, and the gains of the self-improvement loop lasted one or two cycles (\Cref{sec:disc-headline}). The production design (\Cref{app:qa-production-fold}) draws about $3{,}000$ labelled answers per cycle from the stored traffic, so the ratio in deployment depends on the size of the fine-tune it replaces.

% ---------------------------------------------------------------------------- %

\needspace{4\baselineskip}
\section[What the composite catches]{What does the composite catch at inference time that a fine-tune cannot fix?}
\label{app:qa-composite-catches}

A fine-tune updates the generator's weights so that the model is more likely to produce the correct answer on questions it has seen. The fine-tune does not alter the model's behaviour on a specific in-flight query that is hallucinating because of a property the training data did not cover: a fresh trivia question on a recently-changed world fact, an off-topic-but-grounded answer that drifts from the question's intent, a confidently-wrong output where the model has a strong prior toward a misconception. These failure modes are not closeable by parametric updates because they are first-order properties of the model's prior, not properties the training data instances reveal.

The calibrated-probability composite catches each of these failure modes at inference time through a different signal family. Confidently wrong outputs with no passage support trip the early-exit conjunction $(u_{\text{internal}} \ge 0.70) \wedge\allowbreak (p_{\text{ground,max}} \le 0.20)$ and are discarded, so the answer is shown only with the conflicting tag and a warning. Off-topic-but-grounded outputs were meant to be caught by the question-answer match signal $q_{\text{match}}$; on short answers it cannot detect them and acts as a correctness cue instead (\Cref{sec:verifier-qa-rel}), so detecting them is future work. Common-misconception outputs that match the model's prior are the hardest case: the sample-set agreement family (the chain-to-answer entailment $p_{\text{entail}}$, the pairwise similarity $s_{\text{avg}}$ over $M=3$ resampled chains, and the dropout dispersion $u_{\text{dropout}}$ over $K=5$ passes) measures agreement, which a consistently held misconception passes, and the verifier was close to chance on TruthfulQA (\Cref{sec:disc-verifier-scope}). The semantic-entropy signal $h_{\text{norm}}$ was dropped before the main run because its learned weight in the cycle-zero fit was effectively zero while its ten extra samples were the verifier's largest cost. The verifier's value at inference time is separate from what fine-tuning provides: every answer is still checked before it is shown, and answers tagged verified stayed about two-thirds correct at every cycle. The check weakens as the generator changes, however: the share of wrong answers the verifier scored highly rose from $0.208$ for B1 to $0.305$ at C5 (\Cref{tab:headline-steps}). Xiong et al.~\cite{xiong2023can} report that verbalised confidence distinguishes correct from incorrect outputs only weakly (an area under the curve of about $0.6$ for GPT-4) and that no single confidence-elicitation method is consistently best across tasks; the composite therefore aggregates nine signals across four families through per-signal isotonic regression and a regularised logistic boost layer to produce a calibrated correctness probability that does not depend on any one signal.

% ---------------------------------------------------------------------------- %

\needspace{4\baselineskip}
\section[Tier-2 distillation]{What was the zero-shot tier expected to gain from the cycle loop, and did it?}
\label{app:qa-tier2-distillation}

A plain fine-tune produces one checkpoint. The cycle loop produces a sequence of checkpoints trained on stored answers, most of them written by the retrieval-augmented tier and certified by the verifier, so the design expects the fine-tune to fold retrieved knowledge into the generator: later cycles would answer from the zero-shot tier some questions that earlier cycles needed retrieval for, and the retrieval share, and with it the cost per query, would fall.

The run did not show this. The zero-shot share of evaluation queries moved from $48.7$ percent at C0 to $50.2$ percent at C5, and the zero-shot tier's accuracy fell from $0.565$ to $0.519$ (\Cref{tab:tier-trajectory}). Training on retrieval-written chains also taught the generator to cite a context it was not given (\Cref{sec:disc-phantom}). The change in the retrieval share that the run does show comes from the retrieval-utility classifier (\Cref{sec:ruc-ablation}), not from the self-improvement loop.

% ---------------------------------------------------------------------------- %

\needspace{4\baselineskip}
\section[When not to deploy]{When would CAEM not be worth deploying?}
\label{app:qa-when-not}

The architecture is not a free contribution. Each of the six differentiators above corresponds to an architectural component whose cost is non-zero: the verifier ensemble pays forward passes through MiniCheck, the cross-encoder, and the sentence encoder at every query; the labelled cycle-boundary batch requires ongoing labelling effort; the cycle policy requires operational discipline. The architecture is worth deploying exactly when at least one of the six differentiators is binding for the deployment, and is not worth deploying when all six relax simultaneously:

\begin{description}
\item[No abstain class.] Deployments with no requirement for an abstain class.
\item[No scale compounding.] Deployments where memory-direct cost savings do not compound (very low traffic, or a product requirement that every query traverse the same path for uniformity).
\item[Distillation traction absent.] Deployments where the zero-shot tier's per-cycle improvement does not bind because the deployed workload distribution is too narrow for the SIL-driven distillation to gain traction, or too broad for a 3B-class generator to absorb at any cycle horizon.
\item[Manual checkpoint sufficient.] Deployments where a manual previous-checkpoint backup is sufficient and an external snapshot pipeline already solves the catastrophic-forgetting recovery problem.
\item[Labelling abundant.] Deployments with a labelling budget large enough to support frequent direct fine-tuning, where the verifier-as-sample-efficiency-multiplier value collapses because labelling cost is not the binding constraint.
\item[Below parametric ceiling.] Deployments below the parametric-capacity ceiling of the underlying generator, where direct fine-tuning still produces meaningful accuracy gains so the inference-time-error-catching value of the composite is dominated by what fine-tuning would also provide.
\end{description}

A deployment that satisfies all six conditions is one where plain fine-tuning would deliver comparable quality with a simpler implementation. On the panel of this thesis, fine-tuning on gold labels (B6) matched CAEM's exact match and reached a lower composite hallucination metric; CAEM's measured advantages were the status tag and abstention, and a gold-label budget about eleven times smaller (\Cref{ch:evaluation}). A deployment that needs neither a per-answer reliability signal nor a small labelling budget is better served by plain fine-tuning.

\part*{Part II: Design-justification: calibration discipline and cycle-boundary ordering}

The architecture makes several non-obvious design choices in how the calibration fold is constructed, how the cycle boundary is sequenced, and how the deduplication discipline guards the disjointness of the data partition. Part~II answers each as a direct ``is this correct, and why not the simpler alternative'' question.

% ---------------------------------------------------------------------------- %

\needspace{4\baselineskip}
\section[Same calibration fold across cycles]{Why use the same calibration fold across all cycles of the registered horizon rather than rotating folds?}
\label{app:qa-same-fold}

The calibration fold is a fixed $1{,}500$-sample partition ($500$ samples drawn from each of the three training benchmarks) constructed once at pool-builder time and reused at every cycle boundary. The natural alternative would be to rotate the fold across cycles, drawing a fresh $1{,}500$-sample slice each cycle from a larger reservoir. The same-fold protocol is preferred for four reasons.

First, the cycle-boundary recalibration is itself a measurement. The labelled fold is the reference distribution against which the verifier's calibration is read; if the reference distribution varies between cycles, then a cycle-to-cycle change in the fitted curves entangles two sources of variation (the verifier improving under self-improvement, and the reference set drifting) and the trajectory analysis cannot decompose them. A fixed reference fold is the variance-isolation lever that lets the per-cycle isotonic and composite-calibration readings be attributed cleanly to verifier change.

Second, the gold labels on the fold are constant across cycles. The model's verifier signals on a given fold question shift cycle-to-cycle (because the post-fine-tune model produces different token probabilities, dropout variances, chain agreement, and so on), but the question's correct answer does not. Re-fitting the per-signal isotonic curves on the same (variable signals, constant labels) data keeps the calibrator matched to the current model, which matters because, as Ovadia et al.~\cite{ovadia2019trust} found, calibration fitted before a model or data shift degrades after it. Rotating the fold would replace the constant-labels property with a moving-target labelling distribution that confounds the fit.

Third, reusing the fold does not weaken the precision floor of \Cref{thm:purity}: the floor needs new questions to resemble the calibration fold, not one cycle's fold to differ from the last, and the refit at every cycle tracks the drift of the signals.

Fourth, the storage cut-points themselves are held constant across cycles at $\tau_{\text{store}} = 0.60$ and $\tau_{\text{defer}} = 0.45$. With the same fold every cycle and the same cut-points across cycles, any cycle-to-cycle variation in the realised storage precision at the cut-point is attributable to the per-cycle composite refit absorbing distribution drift, rather than to fold-resampling variance or to the cut-points themselves moving.

A defensible alternative for a substantially longer trajectory (say fifty or more cycles) would be a rotating-fold protocol drawn from a $5{,}000$-sample reservoir partitioned into ten overlapping subsets, where each cycle's fold differs but is still an honest sample from the labelled population. The variance-isolation argument weakens at long horizons because the same-fold protocol risks calibration over-fitting to fold-specific noise. The realised six-cycle horizon (the cycle-zero calibration plus cycles one through five, with cycle four aborted under the retention guard and the run stopping after cycle five, as planned after the cycle-four rollback) is short enough that variance isolation dominates, and the rotating-fold extension is registered as a future-work item should the trajectory horizon extend.

% ---------------------------------------------------------------------------- %

\needspace{4\baselineskip}
\section[Production vs research fold]{In production each calibration split would necessarily differ; how is this reconciled with the same-fold research protocol?}
\label{app:qa-production-fold}

In production there is no pre-built calibration partition. The labelled batch at every production cycle boundary is drawn from accumulated serving traffic via stratified sampling (per the production runbook), so each cycle's fold is a fresh sample of about $3{,}000$ labelled queries with content that varies cycle-to-cycle. This is a different operational mode from the research protocol but architecturally compatible with it.

Three points reconcile the two modes. First, the fitting procedure is fold-agnostic. The per-signal isotonic regression is a deterministic monotone fit on whatever (signal, label) pairs are supplied; sklearn's isotonic regression does not care whether the questions were the same as in a previous fit or different. The cycle-zero cut-point selection walks the labelled fold sorted by composite and picks the threshold at the registered precision target; this is run once at cycle zero and the resulting cut-points are held constant across cycles. Both procedures work identically on a fixed fold (research) or a rotating fold (production).

Second, production loses one variance-isolation property and gains a different fidelity property. The research protocol uses a fixed fold to isolate verifier-distribution change as the only source of cycle-to-cycle calibration variation. Production cannot use the fixed fold because the fold material does not exist. Instead, production gains a different property: the calibration fold matches the live deployment distribution exactly, because it is drawn from live deployment traffic. Research's fixed fold is a proxy for deployment (drawn from training-benchmark dev splits); production's rotating fold is the real thing. Each protocol is correct under its own labelling-availability constraint.

Third, the storage and deferral cut-points are fixed in both modes, so a single cycle's traffic anomaly (an unusual spike of out-of-domain queries, for example) cannot move them; it can only change the refitted composite, which the next cycle's labelled batch refits again. The one smoothed quantity is the pre-routing temperature, which blends each cycle's fit with the previous value at weight $\alpha_{\text{ema}} = 0.7$.

The practical consequence: an operator who wants production calibration to behave as close as possible to research-mode discipline can stratify the production sample to be representative of the five-benchmark workload composition and can hold the labelled batch size near the research-mode $1{,}500$-sample point. An operator whose live distribution differs sharply from research can let the production fold reflect the live distribution and accept that the calibration trajectory will diverge from the research trajectory by exactly the distribution-shift magnitude. Both behaviours are correct; the choice is between distribution fidelity and protocol parity.

% ---------------------------------------------------------------------------- %

\needspace{4\baselineskip}
\section[Cycle-boundary ordering]{Why apply the cycle-boundary recalibration before retroactive re-verification rather than after?}
\label{app:qa-cycle-order}

The cycle-boundary order is: self-improvement fine-tune, then score the calibration fold under the now-updated generator, then re-fit the temperature scalar and per-signal isotonic curves and per-benchmark composite boost weights on the freshly-scored fold (the storage cut-points themselves are held constant across cycles), then reload the verifier from the new calibration files, then run retroactive re-verification against the now-recalibrated verifier. The natural alternative would be to run retroactive re-verification first under the previous cycle's calibration and recalibrate afterwards. The current order is the architecturally correct one for two reasons.

First, the post-fine-tune signal distribution falls outside the previous cycle's isotonic fit envelope on a non-trivial fraction of stored episodes. The fine-tune saturates token probability and chain agreement on chains the model trained on; the previous cycle's isotonic curves were fitted on the previous cycle's signal distribution and do not extrapolate cleanly past the upper edge of their training-input range. Running retroactive re-verification under the previous cycle's calibration produces calibrated probabilities on memorised chains that are artefactually low, and the resulting composite drops below the prune threshold for many EM-correct stored entries. An earlier run with the opposite order pruned a large share of the cold-start memory at its first cycle boundary for this reason; the current order refits the calibration before the prune rule reads it.

Second, the corrected order preserves the architectural intent of retroactive re-verification. The re-verification pass is meant to upgrade the stored composite of episodes whose underlying correctness has improved under the new generator and to prune episodes whose correctness has genuinely degraded. The pass cannot do either if the calibration it reads through has not yet been updated to the new model. Under the corrected order, retroactive re-verification reads through fresh per-signal isotonic curves and fresh per-benchmark boost weights at the registered fixed cut-points; an entry whose correctness genuinely improves under the new model maps to a higher calibrated probability and stays in memory with an upgraded $u_{\text{stored}}$, an entry whose correctness genuinely degrades maps to a lower probability and gets pruned. The order is therefore part of what makes the pool-purity floor of \Cref{thm:purity} apply to the re-scored memory.

% ---------------------------------------------------------------------------- %

\needspace{4\baselineskip}
\section[Early-exit disabled at retro-verify]{Why disable the confabulation early-exit during retroactive re-verification when it remains active at query time?}
\label{app:qa-early-exit-retro}

The early-exit rule fires when the internal-confidence composite is high (above $\eta_u = 0.70$) and the top-passage entailment is low (below $\eta_g = 0.20$). At query time this conjunction is the architectural defence against confabulation: a model that is internally confident on a fresh question but has no supporting passage is likely producing a confident hallucination, and the rule discards the answer, which the user sees only with the conflicting tag. At retroactive re-verification time the same conjunction reads differently. The fine-tune just memorised the stored chains, so the internal-confidence composite is automatically near one for every entry the self-improvement loop touched; the first conjunct is no longer informative. Letting the rule fire under this regime collapses the conjunction into ``low top-passage entailment, regardless of confidence,'' which is a much wider catchment than the registered confabulation rule and over-prunes EM-correct stored entries whose chain happened to have thin passage support.

The architectural fix is a query-time-only flag on the rule. The verifier's $\mathrm{verify}$ method takes an $\mathrm{is\_query\_time}$ parameter that defaults to true (preserving query-time inference and deferred-buffer reconsideration behaviour) and is set to false in the retroactive re-verification closure. Under the false setting the early-exit conjunction is short-circuited and the verifier falls through to the standard composite-and-threshold path, which still prunes any entry whose recalibrated $u_{\text{stored}}$ falls below $\tau_{\text{retro}} = 0.50$. Confidently-wrong stored entries are still removed; EM-correct stored entries with thin passage support are no longer falsely pruned. The rule's intended semantics (catching query-time confabulation) is preserved; its inheritance into retroverify-time is corrected.

% ---------------------------------------------------------------------------- %

\needspace{4\baselineskip}
\section[Training vs storage threshold]{Why is the training threshold strictly higher than the storage threshold rather than equal?}
\label{app:qa-tau-train-vs-store}

The architecture registers two different thresholds on the calibrated composite. The storage threshold $\tau_{\text{store}}$ governs admission to the user-facing memory store and to the retrieval tier. The training threshold $\tau_{\text{train}}$ governs admission to the next cycle's self-improvement fine-tune pool; only entries whose composite clears $\tau_{\text{train}}$ enter the gradient signal. The convention is $\tau_{\text{train}} > \tau_{\text{store}}$, with the locked configuration setting $\tau_{\text{store}} = 0.60$ and $\tau_{\text{train}} = 0.70$ (held constant across cycles; the training-pool threshold was lowered from $0.75$ to $0.70$ on the trajectory-launch date to restore a non-empty verified-episode share while preserving the strictly-stricter invariant). The asymmetry is deliberate.

The reason is that the training pool contributes to all subsequent cycles through the generator's weight update, while the storage pool contributes only to the user-facing retrieval tier at the current cycle. A noisy entry in the storage pool affects only the answers served from memory and is re-scored at the next cycle's retroactive re-verification, so the cost of a borderline storage admission is bounded. A noisy entry in the training pool propagates through the gradient signal into every subsequent cycle's generator and can compound the verifier's bias over the trajectory. Natarajan et al.~\cite{natarajan2013learning} show that learning from noisy labels needs the noise to stay below a bound or a loss corrected for the noise rate; CAEM takes the first route, using the stricter threshold to keep the training pool cleaner than the storage pool (about $0.80$ against $0.74$ correct on the stream chunks, \Cref{tab:storage-pool-quality}) rather than correcting the loss.

A defensible alternative would be a single threshold $\tau = \tau_{\text{store}} = \tau_{\text{train}}$ that admits the same entries to both purposes. The simpler alternative loses the asymmetry between storage-class consequences (bounded, single-cycle) and training-class consequences (compounding, multi-cycle), and the training pool is no longer held to a stricter bar than the memory. The two-threshold protocol is the registered defence; the empirical purity of the resulting training pool is reported in \Cref{ch:evaluation}.

% ---------------------------------------------------------------------------- %

\needspace{4\baselineskip}
\section[Three-layer dedup]{Why deduplicate the calibration fold at three layers rather than relying on the dataset's native unique identifiers?}
\label{app:qa-three-layer-dedup}

The deduplication discipline runs at three layers in the pool-builder. The first layer hashes every sample's question text into a sixteen-character content identifier and dedupes within each loaded dataset split. The second layer dedupes within each individual partition slice (calibration, purity, training, evaluation) to catch any residual within-slice duplicates. The third layer cross-checks the calibration and purity slices against the evaluation slice and drops any sample whose content hash collides with an evaluation-slice sample. The natural alternative would be to rely on the dataset's native sample identifiers and assume disjointness by construction.

The dataset's native identifiers are unreliable for two reasons. The lucadiliello FEVER dump reuses integer claim identifiers across the train and dev splits, so a sample with id $13957$ in the train split is content-distinct from a different sample with id $13957$ in the dev split, and a join on native id produces spurious overlap warnings. The TriviaQA dump reuses sfq-prefixed and qw-prefixed and qz-prefixed identifiers within a single split, so two content-identical questions can both appear in the same split with the same id. Relying on native identifiers therefore both over-flags overlap (false positives on different-content same-id pairs) and under-flags overlap (false negatives on same-content different-id pairs). Content-hashing on the question text resolves both regimes by definition: identical questions hash to the same identifier regardless of native id; different questions hash to different identifiers regardless of native id collision.

The three-layer structure is the minimum sufficient cover. A single dataset-level dedup pass leaves intra-slice duplicates untouched. A single intra-slice dedup pass leaves cross-slice content collisions untouched (a calibration question that happens to also appear in the evaluation split). A single cross-slice pass relies on the per-slice slice having been itself dedup-clean; without the within-slice dedup, the cross-slice check sees inflated counts and miscomputes the disjointness invariant. The empirical evidence for the three-layer requirement is reported in the loading-time logs of every cycle: each layer drops some records, and the calibration fold after all three layers has the registered $1{,}500$ questions.
